\documentclass[11pt]{article}
\usepackage{amsmath,amssymb,amsthm,hyperref}
\newtheorem{theorem}{Theorem}
\newtheorem{lemma}{Lemma}
\newtheorem{hypothesis}{Open hypothesis}
\newcommand{\R}{\mathbb R}
\newcommand{\B}{\{0,1\}}
\title{Shared Boolean filters and a weighted partial cover}
\author{Label audit research checkpoint}
\date{24 November 2026}
\begin{document}
\maketitle

\paragraph{Status.}
The reciprocal filter, its positivity, its shared correction basis, and
fixed-weight strictification are kernel checked in
\texttt{HeadComplexity/Results/JointCellFilteredKernels.lean}.
The cover estimate below is a written elementary proof. The proposed
large local domain and its global consequence are conditional on the
explicit decoder hypothesis stated below. No improvement in the
asymptotic unconditional bound is claimed here. All local labels may be
chosen arbitrarily and independently.

\section{Exact filters and arbitrary corrections}

For a finite coordinate set $U$, positive weights $w_i$, and
$d>\sum_{i\in U}w_i$, define
\[
 K_U(w;d)=\sum_{T\subseteq U}
 \frac{(-1)^{|U|-|T|}}{d-\sum_{i\in T}w_i}.
\]
For Boolean $x$, pairing subsets which differ by a missing coordinate gives
\[
 \sum_{T\subseteq U}\frac{(-1)^{|U|-|T|}}
 {d-\sum_{i\in T}w_i x_i}
 =K_U(w;d)\prod_{i\in U}x_i.
\]
Multiplication of the alternating sum by its affine denominator gives
\[
 \left(d-\sum_{i\in U}w_i\right)K_U(w;d)
 =\sum_{i\in U}w_i K_{U\setminus\{i\}}(w;d).
\]
Since $K_\varnothing(w;d)=1/d$, this proves $K_U(w;d)>0$ by induction.
Every denominator is positive on the entire Boolean cube.
The Lean declarations are respectively
\texttt{jointCell\_filtered\_reciprocal\_eq},
\texttt{jointCell\_filtered\_kernel\_recurrence},
\texttt{jointCell\_filtered\_kernel\_pos}, and
\texttt{jointCell\_filtered\_denominator\_pos}, all in namespace
\texttt{HeadComplexity.AxisZeroGluing}.

Let $\mathcal F$ be a finite downward-closed family of coordinate subsets.
Choose one weight vector and one $d$ valid for every member of $\mathcal F$.
For arbitrary real coefficients $c_U$, substitute the preceding identity
for each monomial and collect equal denominators. This gives
\[
 \sum_{U\in\mathcal F}c_U\prod_{i\in U}x_i
 =\sum_{T\in\mathcal F}
 \frac{A_T}{d-\sum_{i\in T}w_i x_i},\qquad
 A_T=\sum_{\substack{U\in\mathcal F\\T\subseteq U}}
 \frac{c_U(-1)^{|U|-|T|}}{K_U(w;d)}.
\]
The coefficients are completely unrestricted. There is one head per
member of $\mathcal F$, not one separate filter per monomial.
This is \texttt{jointCell\_exists\_shared\_reciprocal\_basis}.
For row subsets of size below $c$, write
\[
 D(r,c-1)=\sum_{j=0}^{c-1}\binom rj.
\]
When $1\le c\le r/4$, the ratio between successive lower binomial
coefficients is at most $c/(r-c+1)$. Summing the geometric bound gives
\[
 D(r,c-1)\le\frac{c}{r-2c+1}\binom rc
 \le\frac{2c}{r}\binom rc.
\]
If a correction coefficient depends arbitrarily on both a one-hot bit
coordinate and a one-hot column, a column-linear numerator stores the
column dependence. The row and bit denominator family has size
$(a+1)D(r,c-1)$: each proper row subset occurs alone and with one of the
$a$ bit coordinates. This cost is explicitly charged.

\section{The unresolved decoder interface}

\begin{hypothesis}[Compact filtered decoder]
For given positive integers $c,a,s$, there are $q$ kernels with positive
weights $w_{hi}$, positive bit slopes $\beta_{hb}$, real numerators
$A_{hb}$, constants $D_h$, and $\delta>0$, such that
\[
 \sum_{i=1}^c w_{hi}+\sum_{b=1}^a\beta_{hb}+s\delta<D_h
 \quad(1\le h\le q).
\]
For every word $\sigma\in\{-1,1\}^a$ there are independently chosen
tokens $t_h(\sigma)\in(0,\delta)$ satisfying
\[
 \sigma_b F_b(t(\sigma))>0\quad(1\le b\le a),
 \qquad
 F_b(t)=\sum_{h=1}^q A_{hb}
 \left[K_{[c]}(w_h;D_h-\beta_{hb}-t_h)
       -K_{[c]}(w_h;D_h-\beta_{hb})\right].
\]
The useful regime is bounded $q$, $a$ proportional to $c$, and a validity
range including $s$ proportional to the ambient dimension.
\end{hypothesis}

The independent token vector is allowed at no additional head cost;
requiring a single token shared by the $q$ kernels would be a stronger
hypothesis. Positivity only on the one-hot targets is insufficient.
The displayed budget forces positivity even when all $a$ bit and all
$s$ column coordinates are simultaneously one. In particular the
available token interval must be compact relative to the pole margin.

\section{Conditional large local domain}

Use $r$ row coordinates, $a$ bit coordinates $y$, and $s$ column
coordinates $z$. The target inputs have row support $P$ of size $c$,
$y=e_b$, and $z=e_j$. For every pair $(P,j)$ choose the token vector for
the independently specified $a$-bit label word. For each $(P,h)$ use two
fractions with numerator $\sum_b A_{hb}y_b$ and opposite signs, and
denominators
\[
 D_h-\sum_{i\in P}w_{h,P,i}x_i-\sum_b\beta_{hb}y_b
       -\sum_jt_h(P,j)z_j,
 \qquad
 D_h-\sum_{i\in P}w_{h,P,i}x_i-\sum_b\beta_{hb}y_b.
\]
The decoder budget makes both positive throughout the Boolean cube.
On $(y,z)=(e_b,e_j)$ the row multilinear expansion has support contained
in $P$. Its top coefficient is the corresponding kernel difference.
All lower row coefficients are arbitrary functions of $(b,j)$.
Subtract them using the shared correction family from the preceding
section. On a row support $P$ of size $c$, exactly its own top monomial
survives, so the score is $F_b(t(P,j))$ with the required strict sign.
On every smaller row support it is zero before strictification.
The number of heads is
\[
 H_{\rm local}\le 2q\binom rc+(a+1)D(r,c-1).
\]
Thus $\binom rc as$ arbitrary labels are handled by one score, with
an explicit correction charge rather than assumed cancellation of
independent label coefficients.

The bit axis $y=0$ vanishes: the main numerators vanish and each
correction filter contains a bit coordinate. The column axis $z=0$
vanishes: paired main denominators agree, and correction numerators
vanish. These cancellations survive addition of the same affine
offset to every denominator in this local module. Pairing subsets in
the bit filter remains valid even when that common offset depends on
other coordinates.

\paragraph{Strict slopes.}
The displayed raw denominators have weak negative slopes. Add
$\varepsilon(c+2-|x|-|y|-|z|)$ to every denominator, with the same
positive sufficiently small $\varepsilon$. Every slope becomes
strictly negative, target denominators are unchanged, and positivity
on the finite Boolean cube persists. Both axes remain exactly zero.
The zero scores on smaller row supports may cease to be zero, but
converge uniformly to zero as $\varepsilon$ tends to zero. This is the
quantity needed to preserve previously fixed strict signs.
The single-head fixed-weight assertion, including original-model atom
admissibility and full-cube positivity, is kernel checked as
\texttt{jointCell\_exists\_atom\_on\_weight\_slice}. Its explicit
choice of $\varepsilon$ is per head; a finite module uses the minimum
of the finitely many positive margin bounds, and continuity also
allows the sign-preservation bounds below.

\section{Weighted partial covering of the row cube}

The following estimate avoids a growing cover-density charge by
paying directly for uncovered rows. It is not a claim that exact
$c$-deletion stars alone have constant covering density.

Let $m$ be the row cube dimension, $1\le c\le m/16$, and suppose a
$c$-deletion star with center of size $r\ge m/4$ costs at most
$K\binom rc$ heads, where
\[
 K=2q+\frac{8(a+1)c}{m}.
\]
This follows from the local count above. Let $B=4as$ be the cost of a
singleton row fallback. A fallback is constructed entry by entry:
for each $(b,j)$, the two-coordinate reciprocal filter on
$\{y_b,z_j\}$ has four heads, a positive target value, and zero on
both axes. Scale its numerator by the desired sign and by the inverse
positive target value. It remains axis-zero after common affine
denominator offsets. Strictification and row penalties therefore work
for this fallback as well.

For a target row of weight $k$, eligible centers have size $k+c$.
Their number above that row is $\binom{m-k}{c}$. Independently select
each center on this layer with probability
\[
 p_k=\min\left\{1,\frac{t}{\binom{m-k}{c}}\right\},\qquad t\ge1.
\]
Only layers with $k+c\ge m/4$ and $k+c\le m$ are eligible.
If $p_k=1$, no target row on the layer is uncovered. Otherwise its
uncovered probability is at most $e^{-t}$. For selected-center cost,
the incidence identity gives
\[
 \frac{\binom m{k+c}\binom{k+c}c}{\binom{m-k}c}
 =\binom mk.
\]
Consequently expected selected-center cost, including every overlap,
is at most $Kt2^m$. Expected fallback cost on eligible layers is at
most $Be^{-t}2^m$.

Excluded low layers satisfy $k<m/4$, and excluded high layers satisfy
$k>m-c\ge3m/4$. If $X$ is binomial with parameters $(m,1/2)$, Markov's
inequality applied to $3^{-X}$ gives
\[
 \Pr(X\le m/4)\le3^{m/4}(2/3)^m=e^{-\kappa m},
 \qquad \kappa=\tfrac34\log3-\log2>0.
\]
By symmetry the high tail has the same bound. The excluded rows cost
at most $2B e^{-\kappa m}2^m$ heads. Hence some deterministic
selection and its fallback rows have total cost at most
\[
 2^m\bigl(Kt+Be^{-t}+2Be^{-\kappa m}\bigr).
\]
Taking $t=\log B$ and $m\ge2\log B/\kappa$ yields the useful bound
\[
 H_{\rm row}\le2^m\bigl(K\log(4as)+3\bigr).
\]
The logarithm depends on the fallback table cost, not on the size
$\binom rc$ of an edge. This is the source of the improvement over
a naive union bound demanding that every row be covered by stars.

\section{Strict signs, penalties, and orientation}

Fix a bit-column star pair. Process all selected row stars and
singleton fallbacks in decreasing order of their \emph{target row
weight}. A star centered at $X$ has target weight $|X|-c$, whereas a
singleton centered at $Y$ has target weight $|Y|$. This order is
essential; ordering by center weight alone is insufficient when
stars and fallbacks have different radii.

At a new star centered at $X$, an earlier protected input inside its
downset has deletion count at most $c$. If the count equals $c$, it
is a new target too and the required label agrees. If it is below
$c$, the unperturbed local score is zero. First choose a positive
scale large enough to dominate the existing finite score on all new
targets. Then choose the common strictification parameter small
enough that the zero-limit scores on earlier inside targets preserve
their strict signs. These choices are possible because all sets are
finite and the new target margins are strictly positive.

Finally add a common nonnegative penalty
$T\sum_{i\notin X}u_i$ to every local denominator, in ambient row
coordinates $u$. It vanishes on the downset of $X$. On each earlier
protected input outside that downset it tends to infinity as
$T\to\infty$, so the new module tends to zero there. Choose $T$
large enough to preserve these finitely many earlier strict signs.
No denominator loses positivity. All axis cancellations survive
because the penalty is common to every denominator in the module.
The same argument handles a singleton fallback with radius zero.

The assembled row module thus has every prescribed sign on the full
row cube times the chosen bit-column star pair, and vanishes exactly
on its bit or column axis for all row inputs. Apply the inherited
descending-center two-star gluing to the bit and column blocks.
Inside a new bit-column downset, earlier inputs are either on one
axis, where the module is exactly zero, or in the new star pair,
where labels agree. Outside the downset a common positive penalty
suppresses interference as before.

In local deletion coordinates the slopes are strictly negative.
Complementing coordinates inside a center turns these slopes
strictly positive in ambient original-model coordinates. Outside
coordinates receive strictly positive filler slopes and
nonnegative selection penalties. The row stage may temporarily have
positive row slopes and negative bit-column slopes; the final
bit-column coordinate complementation makes every ambient slope
positive. Coordinate complementation preserves full-cube positivity,
and all numerators stay affine. Thus the construction would use the
original admissibility conditions, without modifying the model.

\section{Conditional asymptotic consequence and the bottleneck}

Use the existing constant-factor one-deletion covers for a bit block
of dimension $A$ and a column block of dimension $S$, and the row
partial cover just proved. The top vertices omitted by one-deletion
stars can be paid separately; their relative proportion is
exponentially small in $A$ or $S$. A decoder uniform for $a=A,s=S$
restricts to smaller star sizes by extending each label word
arbitrarily. Let $n=m+A+S$, with $S$ proportional to $m$ and
$A,c$ proportional to $\sqrt m$. If the compact decoder exists with
bounded $q$, then $K=O(1)$ and the global count is
\[
 W(n)=O\left(\frac{2^n\log n}{n^{3/2}}\right).
\]
This would imply the requested fixed-power bound for every fixed
$\delta<1/2$. The local input count contains
$\binom rc=\exp(\Theta(c\log(r/c)))$, so it passes the independent-cell
size test. The missing lemma is the compact decoder, not the cost of
independent low-degree corrections or the row cover density.

\paragraph{Finite evidence, not a theorem.}
The script \texttt{JointCellFilteredKernelSearch.py} searches ratios
of two positive filtered kernels. The stored file
\texttt{joint-cell-filtered-c8-candidate.json} contains one pair of
eight-weight kernels with three numerically observed derivative
sign changes, reevaluated at 70 decimal digits. This does not prove
the number of extrema, an exponential number of sign words, the
compact-interval requirement, or any decoder hypothesis. A
single positive kernel, or a positive sum with no signed
competition, cannot produce the required arbitrary words.

\end{document}
